%!TEX root = ./../../thesis-trier.tex

\IN{factor (semi-)automaton} \label{def:factor_aut}
Let $\semiautomat{A}{Q}{\Sigma}{\delta}$. Suppose $\sim$ is a $\TS{\mathcal A}$-congruence. 
Equivalently, for every $p,q \in Q$ and $x \in \Sigma$,
we have
\begin{equation}\label{prelim:eqn:TS-congruence}
	p \sim q \Rightarrow \applytrans{\delta}{p}{x} \sim \applytrans{\delta}{q}{x}.
\end{equation}
The \emph{quotient automaton} is $\semiautomat{A / {\sim}}{Q / {\sim}}{\Sigma}{\delta_{\sim}}$ 
with the well-defined transition function $\applytrans{\delta_{\sim}}{[q]_{\sim}}{x} = [\applytrans{\delta}{q}{x}]_{\sim}$
for all $x \in \Sigma$ and~$q \in Q$.

The next two statements connect $\TS{\mathcal A}$-congruences, quotient automata, the kernel relation, and morphisms
in the following way: $\TS{\mathcal A}$-congruences
correspond to surjective morphisms onto quotient automata
via the canonical morphism (\cref{prelim:prop:congruence_induces_morphism}), and surjective morphisms
correspond to $\TS{\mathcal A}$-congruences via the kernel relation (\cref{prelim:thm:homtheorem}).

\IN{canonical morphism}
In view of~\cref{prelim:prop:congruence_induces_morphism}, in this context we also call $\pi_{\sim}$
the \emph{canonical morphism (for~$\sim$)} 
from $\mathcal A$ onto $\mathcal A / {\sim}$. 

\begin{theorem}
	\label{prelim:prop:congruence_induces_morphism}
	Let $\semiautomat{A}{Q}{\Sigma}{\delta}$
	with $\TS{\mathcal A}$-congruence $\sim$ on $\mathcal A$.
    Then the canonical projection map $\pi_{\sim} \colon Q \to Q / {\sim}$ from \cref{def:canprojmap}
    is a surjective morphism whose kernel relation equals~$\sim$.
\end{theorem}
\begin{proof}
	Clearly, $\pi_{\sim}$ is surjective and its kernel relation equals $\sim$.
	It is a morphism since by definition of $\mathcal A / {\sim}$
	we have $\delta_{\sim}(\pi_{\sim}(q), x) = \pi_{\sim}(\delta(q, x))$
	for $q \in Q$ and $x \in \Sigma$.
\end{proof}

\begin{theorem} 
	\label{prelim:thm:homtheorem}
	Let $\varphi \colon Q \to S$ be a morphism
	between the automata $\semiautomat{A}{Q}{\Sigma}{\delta}$
	and $\semiautomat{B}{S}{\Sigma}{\mu}$.
	Then the following is true:
	\begin{enumerate} 
	\item \label{prelim:thm:homtheorem:1} $\applyfunc{Q}{\varphi}$ is an upper set in $\mathcal B$,
	
	\item \label{prelim:thm:homtheorem:2} the kernel relation $\sim_{\varphi}$
	is a $\TS{\mathcal A}$-congruence,
	
	\item \label{prelim:thm:homtheorem:3} $[q]_{\sim_{\varphi}} \mapsto \applyfunc{q}{\varphi}$, $q \in Q$,
	is a well-defined isomorphism between
	$\mathcal A / {\sim_{\varphi}}$ and~$\mathcal B_{|\applyfunc{Q}{\varphi}}$.
	\end{enumerate}
\end{theorem}
\begin{proof}
\ref{prelim:thm:homtheorem:1}: This follows from \crefitem{prelim:automata:morpisms:reachability:prop:images}{prelim:automata:morpisms:reachability:prop:images:preservation}.

\ref{prelim:thm:homtheorem:2}: Suppose $p \sim_{\varphi} q$ for $p,q \in Q$.
 Then we have $\varphi(p) = \varphi(q)$.
 Since $\varphi$ is a morphism,
 we find $\varphi(\delta(p, x)) = \mu(\varphi(p), x) = \mu(\varphi(q), x) = \varphi(\delta(q, x))$,
 and so $\delta(p, x) \sim_{\varphi} \delta(q, x)$, i.e., $\sim_{\varphi}$
 is a $\TS{\mathcal A}$-congruence.
 
\ref{prelim:thm:homtheorem:3}: Well-definedness follows from the definition of the kernel relation.
 Now $\delta_{\sim_{\varphi}}([q]_{\sim_{\varphi}}, x) = [\delta(q, x)]_{\sim_{\varphi}}$ 
 is mapped to $\varphi(\delta(q, x))$
 and $[q]_{\sim_{\varphi}}$ is mapped to $\varphi(q)$.
 Since $\mu(\varphi(q), x) = \varphi(\delta(q, x))$,
 we obtain that the mapping is a morphism
 as we arrive at the same result when applying first a letter and then the mapping,
 or when first applying the mapping and then a letter.
 The mapping clearly maps onto $\mathcal B_{|\varphi(Q)}$.
\end{proof}

A similar statement as \cref{prelim:thm:homtheorem} in the theory of groups is often called the 
homomorphism theorem, for example, see~\cite[Sect.~1.2.5]{Kurzweil04}.
\Cref{prelim:thm:homtheorem} states that morphisms composed along paths in the following diagram
ending at the same point are equal.
\[
\begin{tikzcd}[row sep=1.3cm, column sep=2.5cm]
	Q \arrow{r}{q \mapsto [q]_{\sim_{\varphi}}}
	\arrow[swap,
	start anchor={[shift={(-3pt,3pt)}]south east}, % fine tune start of arrow
	end anchor={[shift={(2pt,-4pt)}]north west} % fine tune end of arrow
	]{dr}{\varphi}
	& Q / {\sim_{\varphi}} \arrow{d}{[q]_{\sim_{\varphi}} \mapsto \varphi(q)} \\
	& \varphi(Q) 
\end{tikzcd}
\]

\begin{corollary}	\label{prelim:cor:homtheorem}
	Let $\varphi \colon Q \to S$ be a morphism
	between the automata $\semiautomat{A}{Q}{\Sigma}{\delta}$
	and $\semiautomat{B}{S}{\Sigma}{\mu}$ with $Q \ne \emptyset$.
	If $\mathcal B$ is strongly connected, 
	then $\applyfunc{Q}{\varphi} = S$, i.e., $\varphi$ is surjective. So $\mathcal A / {\sim_{\varphi}} \cong \mathcal B$. 
	In particular, for $\mathcal A = \mathcal B$,
	by finiteness, $\varphi$ is an isomorphism.	
\end{corollary}

\begin{corollary}
  Suppose $\semiautomat{A}{Q}{\Sigma}{\delta}$ admits
  no non-trivial $\TS{\mathcal A}$-congruence.
  Then every morphism from $\mathcal A$ into another automaton
  is either injective or constant.  
\end{corollary}
\begin{proof}
  For a morphism into another automaton, its kernel relation
  is a $\TS{\mathcal A}$-congruence by \crefitem{prelim:thm:homtheorem}{prelim:thm:homtheorem:2}.
  So it either equals the universal relation, and so the morphism
  is a constant, or the equality relation, and so the morphism is injective.  
\end{proof}

Next, we show a condition to obtain a morphism from two morphisms
from the same automaton. The proof is straightforward and omitted.

\begin{theorem}
	\label{prelim:thm:homsfactorthrough}
	Let $\varphi \colon Q \to S$
	and $\psi \colon Q \to T$
	be morphisms
	from $\semiautomat{A}{Q}{\Sigma}{\delta}$
	to $\semiautomat{B}{S}{\Sigma}{\mu}$
	and $\semiautomat{C}{T}{\Sigma}{\xi}$, respectively.
	Suppose $\sim_{\varphi}$ refines $\sim_{\psi}$.
	Then $\varphi(q) \mapsto \psi(q)$
	is a well-defined surjective morphisms
	from $\mathcal B_{|\varphi(Q)}$
	onto $\mathcal C_{|\psi(Q)}$.
\end{theorem}

By \crefitem{prelim:thm:homtheorem}{prelim:thm:homtheorem:3},
$\mathcal A / {\sim_{\varphi}} \cong \mathcal B_{|\varphi(Q)}$
and $\mathcal A / {\sim_{\psi}} \cong \mathcal C_{|\psi(Q)}$
in \cref{prelim:thm:homsfactorthrough}. 
In the same way as for \cref{prelim:thm:homtheorem}, the situation from \cref{prelim:thm:homsfactorthrough} is illustrated in the following diagram.
\[
\begin{tikzcd}[column sep=1.3cm, row sep=1.3cm]
	& Q \arrow[dl, swap, "\varphi"] \arrow[dr, "\psi"] & \\
	S \arrow{rr}{\varphi(q) \mapsto \varphi(q)} & & T
\end{tikzcd}
\]

\begin{corollary}
 Let $\semiautomat{A}{Q}{\Sigma}{\delta}$ and $\semiautomat{B}{S}{\Sigma}{\mu}$
 with $\sim_{\mathcal A}$ and $\sim_{\mathcal B}$
 being $\TS{\mathcal A}$- and $\TS{\mathcal B}$-congruences, respectively.
 If $\sim_{\mathcal A}$ refines $\sim_{\mathcal B}$,
 then $[q]_{\sim_{\mathcal A}} \mapsto [q]_{\sim_{\mathcal B}}$
 is a morphism from $\mathcal A / {\sim_{\mathcal A}}$
 into $\mathcal B  / {\sim_{\mathcal B}}$.	
\end{corollary}